\documentclass[10pt,a4paper]{amsart}
\usepackage[margin=19mm]{geometry}
\usepackage{amsmath,amssymb,mathtools}
\usepackage[scr=rsfs,cal=euler]{mathalfa}
\usepackage{xfrac}
\usepackage[unicode,hidelinks,bookmarks=false]{hyperref}
\newcommand{\ie}{i.e.\ }
\newcommand{\cf}{cf.\ }
\newcommand{\resp}{respectively\ }
\newcommand{\Subsection}{\subsection}
\providecommand{\nobreakdash}{\nobreak\hbox{-}}
\setlength{\emergencystretch}{2em}
\multlinegap0pt
\usepackage{bm}
\usepackage{bbm}
\usepackage{dsfont}
\usepackage{xspace}
\usepackage{cancel}
\usepackage{mathtools}
\usepackage{cleveref}
\usepackage[shortlabels]{enumitem}
\usepackage{amsthm}
\makeatletter
\@ifundefined{lemma}{\newtheorem{lemma}{Lemma}}{}
\@ifundefined{proposition}{\newtheorem{proposition}{Proposition}}{}
\@ifundefined{theorem}{\newtheorem{theorem}{Theorem}}{}
\makeatother
\DeclareMathOperator{\Comp}{Comp}
\DeclareMathOperator{\Fact}{Fact}
\newcommand{\Hilbert}{\mathcal{H}}
\newcommand{\la}{\langle}
\newcommand{\rand}{\partial}
\newcommand{\ra}{\rangle}
\newcommand{\xto}[1]{\stackrel{#1}{\to}}
\renewcommand{\phi}{\varphi}
\title[Boundary values via reproducing kernels: The Julia-Carath\'e]{Boundary values via reproducing kernels: The Julia-Carath\'eodory theorem}
\author{Frej Dahlin}
\date{Source: arXiv:2412.13901v2 (CC BY 4.0)}
\begin{document}
\maketitle

\noindent\emph{Source and licence.} Boundary values via reproducing kernels: The Julia-Carath\'eodory theorem. Version arXiv:2412.13901v2 (CC BY 4.0), distributed under the Creative Commons Attribution 4.0 International licence, as recorded in the arXiv licence metadata for this exact version and shown on the abstract page https://arxiv.org/abs/2412.13901v2. Full rights notice: licenses/arxiv-2412.13901-CC-BY-4.0.txt.\par\smallskip

\noindent\textit{Editorial scope.} Counted unit: the Theorem whose source label is \texttt{theorem:main} in the article (source0.tex lines 866--886, source label \texttt{theorem:main}), delivered together with the article's own complete original proof of it (source0.tex lines 887--943, one \texttt{proof} environment of 57 lines). No mathematics is written, repaired or re-derived: statement and proof are the article's own text line for line. Required context retained uncounted: boundary-to-diagonal (lines 244--256); compact-boundary (lines 244--256); constants-inclusion (lines 244--256); isolated-singularity (lines 244--256); proposition:weighted-composition (lines 445--453); lemma:boundary-explosion (lines 567--570); lemma:growth-restriction (lines 608--616). Every definition, display and label of those lines is retained unchanged. Omitted from this package and not redistributed: 1--243, 257--444, 454--566, 571--607, 617--865, 944--1268. Nothing inside a retained excerpt is omitted or altered: every retained body is delivered exactly as the article's lines are written, so no omission receipt of any kind accompanies this package. Citations: the retained excerpts carry no citation command at all, so no bibliography entry of the article is transcribed and no reference is redistributed. Reference adaptation: none was needed and none was made; every reference inside the delivered unit resolves inside it, so no unresolved reference or citation is produced. Printed slips kept literal: no printed word, symbol, hypothesis or number of the counted statement or of its proof was altered. Layout and class: the article's own class is \texttt{article} with packages including ifdraft, geometry, microtype, amsthm, amsmath, amssymb, mathtools, enumitem, biblatex, cleveref; the wrapper is the lane's pinned \texttt{amsart} editorial wrapper at 10pt on A4, which restates the theorem-like environments the retained text needs and adds the article's own macro definitions that the retained text needs (\texttt{\textbackslash Comp}, \texttt{\textbackslash Fact}, \texttt{\textbackslash Hilbert}, \texttt{\textbackslash la}, \texttt{\textbackslash phi}, \texttt{\textbackslash ra}, \texttt{\textbackslash rand}, \texttt{\textbackslash xto}), with extra wrapper packages (bm, bbm, dsfont, xspace, cancel, mathtools, cleveref, enumitem). The delivered numbering is the unit's own. Assets: no retained span contains a figure environment, a graphics-inclusion command, a PDF-page inclusion, a table or a TikZ picture; no image file is copied into this package, the package declares no asset dependency and no image file is redistributed with it. Licence: version arXiv:2412.13901v2 of this article is distributed under the Creative Commons Attribution 4.0 International licence, as recorded in the arXiv licence metadata for this exact version and shown on the abstract page https://arxiv.org/abs/2412.13901v2. A full rights notice is delivered at licenses/arxiv-2412.13901-CC-BY-4.0.txt. Characters: every retained excerpt is pure ASCII, so the delivered engine, XeLaTeX with its default Latin Modern text font, reports no missing character.
\begin{enumerate}[label=(\Alph*)]
\item\label{boundary-to-diagonal}
For every $\xi\in\rand_kX$ there exists a constant $a_\xi$ such that
$|k_\xi(x)| \leq a_\xi k(x,x), \forall x\in X.$
\item\label{constants-inclusion}
There is a constant $b>0$ such that $|k(x,y)|>b$ for all $x,y\in X$.
\item\label{isolated-singularity}
Let $\xi\in\rand_kX$ and $(x_n)$ be a sequence in $X$. If $|k_\xi(x_n)|\to\infty$,
then $x_n\xto{k}\xi$.
\item\label{compact-boundary}
If for a sequence $(x_n)$ in $X$ $k(x_n,x_n)\to\infty$ as $n\to\infty$,
then there is a subsequence such that $x_{n_j}\xto{k}\xi\in\rand_kX$.
\end{enumerate}

\begin{proposition}\label{proposition:weighted-composition}
Let $\phi\in\Fact(k,t)$, then for every $f\in\Hilbert(\frac{k}{t\circ\phi})$ the
pair $(f,\phi)$ induces a bounded weighted composition operator from $\Hilbert(t)$
into $\Hilbert(k)$ of norm at most $\|f\|$. That is, the linear map defined by
$(W_{f,\phi})g(x) = f(x)g(\phi(x))$ is bounded from $\Hilbert(t)$ into $\Hilbert(k)$.
\\ In particular, if the constant function $x\mapsto1$ belongs to $\Hilbert(\frac{k}{t\circ\phi})$,
then $\phi\in\Comp(t, k)$ and
if $x\mapsto1\in\Hilbert(t\circ\phi)$, then $\Hilbert(\frac{k}{t\circ\phi})\subset\Hilbert(k)$.
\end{proposition}

\begin{lemma}\label{lemma:boundary-explosion}
Let $x_n\xto{k}\xi\in\rand_kX$, then $\|k_{x_n}\|^2 = k(x_n,x_n)\to\infty$ as
$n\to\infty$.
\end{lemma}

\begin{lemma}\label{lemma:growth-restriction}
If $x_n\xto{k}\xi\in\rand_kX$, then the normalized kernel functions
$\hat k_x = \frac{k_x}{k(x,x)^\frac12}$ tend weakly to $0$ along this sequence,
that is, for every $f\in\Hilbert(k)$ \begin{equation}
	\lim_{n\to\infty}\la f,\hat k_{x_n}\ra
	= \lim_{n\to\infty}\frac{f(x_n)}{k(x_n,x_n)^\frac12}
	= 0.
\end{equation}
\end{lemma}

\begin{theorem}\label{theorem:main}
Let $k$ and $t$ be reproducing kernels on nonempty sets $X$ and $Y$, respectively.
Suppose that $t$ satisfies \ref{boundary-to-diagonal}-\ref{compact-boundary}.
For $\phi\in\Fact(k,t)$ and $\xi\in\rand_kX$, suppose that $\Gamma_k(M,\xi)$
approaches $\xi$ for some $M>0$, then the following are equivalent:
\begin{enumerate}[label=(\roman*)]
\item\label{item:liminf-condition} \[
	\liminf_{x\xto{k}\xi}\frac{k(x,x)}{t\circ\phi(x,x)} = c < \infty.
\]
\item\label{item:kernel-quotient}
There is a $\lambda\in\rand_tY$ such that the function defined by \[
	q_\xi(x) = \frac{k_\xi(x)}{t_\lambda(\phi(x))}
\]
belongs to $\Hilbert(\frac{k}{t\circ\phi})$.
\item\label{item:uniform-boundary-values}
All functions in $\Hilbert(\frac{k}{t\circ\phi})$ have a $\Gamma_k$-limit at $\xi$.
\end{enumerate}
Furthermore, if any of the above hold, $\phi$ has $E_k$-limit $\lambda$ at $\xi$,
and we have the estimate $0<\|q_\xi\|^2\leq c \leq (Ma_\lambda\|q_\xi\|)^2$ for
every $M$ such that $\Gamma_k(M,\xi)$ approaches $\xi$.
\end{theorem}

\begin{proof}
Fix the notation $q=\frac{k}{t\circ{\phi}}$.

Suppose that \ref{item:liminf-condition} then by assumption there exists a
sequence $x_n\xto{k}\xi$ such that $\|q_{x_n}\|^2 = q(x_n,x_n)\to c$. By
\cref{lemma:boundary-explosion}, $k(x_n,x_n)\to\infty$ as $n\to\infty$,
therefore $t\circ\phi(x_n,x_n)\to\infty$ as $n\to\infty$. Hence, by weak
compactness and \ref{compact-boundary}, passing to a subsequence if
necessary, there is a function $q_\xi\in\Hilbert(q)$ such that $q_{x_n}\xto{w}q_\xi$,
as well as a $\lambda\in\rand_tY$ such that $\phi(x_n)\xto{t}\lambda$. We
compute the function $q_\xi$ as \begin{equation}\label{equation:qxi-computation}
	q_\xi(y)
	= \la q_\xi, q_y\ra
	= \lim_{n\to\infty}\la q_{x_n}, q_y\ra
	= \lim_{n\to\infty}\frac{k_{x_n}(y)}{t(\phi(y),\phi(x_n))}
	= \frac{k_\xi(y)}{t_\lambda(\phi(y))}.
\end{equation}

Suppose that \ref{item:kernel-quotient}, and let $x_n\xto{E_k}\xi$.
\cref{proposition:weighted-composition} together with \ref{constants-inclusion}
says that the function $q_\xi$ belongs to $\Hilbert(k)$. Therefore, by \cref{lemma:growth-restriction} \[
	\frac{q_\xi(x_n)}{k(x_n,x_n)}
	= \frac{k_\xi(x_n)}{t_\lambda(\phi(x_n))k(x_n,x_n)^\frac12}
	\to 0.
\]
But since $x_n\in E_k(M,\xi)$ for some fixed $M>0$, that is,
$\frac{|k_\xi(x_n)|}{k(x_n,x_n)^\frac12}\geq\frac1M$, we have
$|t_\lambda(\phi(x_n))| \to \infty$, hence $\phi(x_n)\xto{t}\lambda$
because of \ref{isolated-singularity}.
\\ Now let $x_n\to\xi$ within $\Gamma_k(M,\xi)$ for some fixed $M>0$,
\ref{item:uniform-boundary-values} is equivalent to showing that the
functions $q_{x_n}$ tend weakly to $q_\xi$ as $n\to\infty$. In the previous
paragraph we established pointwise convergence, since the linear span of
$\{q_x:x\in X\}$ is dense in $\Hilbert(q)$ it is then enough to prove that the
sequence $(q_{x_n})$ is bounded in norm. To this end, we estimate as follows,
for any $x\in\Gamma_k(M,\xi)$ we have \begin{equation}\label{equation:c-norm-estimate}
	\|q_{x}\|^2
	= \frac{k(x,x)}{t\circ\phi(x,x)}
	\leq \frac{|k_\xi(x)|}{t\circ\phi(x,x)}
	\leq Ma_\lambda\frac{|k_\xi(x)|}{|t_\lambda(\phi(x))|}
	\leq Ma_\lambda\|q_\xi\|\|q_x\|,
\end{equation}
where the three inequalities are given by $x\in\Gamma_k(M,\xi)$,
\ref{boundary-to-diagonal}, and Cauchy-Schwarz, respectively.

Suppose that \ref{item:uniform-boundary-values}, as stated previously,
this is equivalent to the sequence $(q_{x_n})$ tending weakly to $q_\xi$
whenever $x_n\xto{\Gamma_k}\xi$. Therefore the implication to \ref{item:liminf-condition}
is given by the principle of uniform boundedness.

The computation \eqref{equation:qxi-computation} reveals that $c\geq\|q_\xi\|^2$, which in
turn satisfies $\|q_\xi\|>0$, for otherwise $q_\xi$ is the zero function so $k_\xi$ is
as well, contradicting the fact that $k_\xi\notin\Hilbert(k)$.
In the proof of $(ii)\implies(iii)$ it was established that $\phi$ has $E_k$-limit
$\lambda$ at $\xi$.
The final desired estimate was given by \eqref{equation:c-norm-estimate}.
\end{proof}

\end{document}
